% =============================================================================
% 3 TRABALHOS RELACIONADOS — segunda versão (18/11/2026), encorpada.
% Organização por tema, e não por autor (skill estruturacao-introducao-
% logica). Cada seção fecha com o que a linha deixa em aberto; a Seção 3.5
% junta as lacunas. Todo número de outro trabalho vem da fonte citada,
% conferida na fonte primária em 18/11/2026 (ver referencias.bib).
% =============================================================================
\chapter{Trabalhos relacionados}
\label{cap:trabalhos-relacionados}

Esta seção situa o trabalho em quatro linhas: os \textit{benchmarks}
de \textit{text-to-SQL}, os estudos sobre o contexto fornecido ao modelo,
as camadas semânticas usadas com modelos de linguagem e a governança de
agentes que consultam bancos de dados. Cada seção termina com o que a
linha deixa em aberto, e a \autoref{sec:lacuna} reúne essas lacunas na
questão de pesquisa do trabalho.

% -----------------------------------------------------------------------------
\section{Benchmarks de \textit{text-to-SQL}}
\label{sec:benchmarks}

O Spider estabeleceu o cenário de avaliação entre domínios, com 10.181
perguntas e 5.693 consultas sobre 200 bancos de 138 domínios, e bancos
diferentes no treino e no teste: o modelo precisa generalizar para
esquemas que não viu \cite{yu2018spider}. Os bancos vieram de material
didático, como disciplinas, tutoriais e livros de SQL, de modelos de
esquema populados com dados gerados e de tabelas da Wikipédia
reorganizadas. Na versão original, a métrica principal era a
correspondência exata da consulta. São bancos pequenos e organizados,
adequados para medir a tradução da pergunta, mas distantes de uma base
operacional.

O BIRD aproximou a avaliação de dados reais \cite{li2023bird}. Seus 95
bancos, com 33,4~GB no total, cobrem 37 domínios profissionais e vêm de
fontes abertas: cerca de um terço do Kaggle, metade de um repositório
acadêmico de bancos relacionais e o restante de tabelas abertas
organizadas pelos autores. Os valores são ``sujos'', como salários
armazenados em texto com símbolo de moeda, e a avaliação usa a acurácia
de execução. Cada banco traz um arquivo de descrição das colunas, e cada
pergunta traz uma evidência de conhecimento externo escrita por
anotadores, com fórmulas, sinônimos ou valores. Antes dele, o KaggleDBQA
já usava bancos reais publicados na web e mostrava que a documentação do
banco melhorava os analisadores da época \cite{lee2021kaggledbqa}.

\textit{Benchmarks} mais recentes tratam do cenário corporativo. O Spider
2.0 reúne 632 tarefas de fluxos de trabalho reais, sobre bancos com
centenas de colunas em sistemas como BigQuery, Snowflake e PostgreSQL. Um
agente com o modelo o1-preview resolveu 21,3\% delas, contra 91,2\% no
Spider original e 73,0\% no BIRD \cite{lei2025spider2}. O BEAVER reúne
9.128 pares de pergunta e consulta sobre 812 tabelas de \textit{data
warehouses} privados, extraídos de registros reais de consultas e
complementados por consultas sintetizadas; agentes com o GPT-5.2 chegaram
a cerca de 10,8\% de acurácia, e a cerca de 30,1\% quando receberam
anotações-oráculo para as subtarefas \cite{chen2024beaver}.

Em português, a linha começou com a tradução do Spider e o treinamento de
modelos bilíngues \cite{jose2021mrat}. Avaliações mais recentes compararam
17 modelos de linguagem no Spider em português e em inglês, com o esquema
apresentado como texto ou como definição das tabelas
\cite{carvalho2025benchmarking}. Com dados públicos brasileiros, o
JABUTI-SQL reúne 69 pares sobre dados reais do DATASUS, com as
inconsistências típicas de bases governamentais \cite{botelho2026jabuti}.
Em um \textit{benchmark} unificado em português, seis modelos passaram de
70\% de acurácia de execução no Spider, mas o melhor resultado sobre a
base do DATASUS foi 27,5\% \cite{fonseca2026avaliacao}.

Esses trabalhos mostram que a acurácia cai quando a base se aproxima da
realidade, mas medem modelos e sistemas, e não o efeito do tipo de
contexto. Nos que fornecem conhecimento, como o BIRD, ele vem pronto e por
pergunta, como uma dica, e não como um conjunto de regras da base mantido
pela organização. Nenhum deles varia, sobre uma mesma base comercial, o
tipo de informação que o modelo recebe.

% -----------------------------------------------------------------------------
\section{Engenharia de prompt e contexto em \textit{text-to-SQL}}
\label{sec:contexto-externo}

Com modelos de linguagem, a tradução de perguntas em SQL passou a ser
feita por meio do prompt, e a forma de descrever o banco tornou-se objeto
de estudo. Os primeiros trabalhos compararam formatos de esquema com o
modelo Codex: o texto de criação das tabelas acompanhado de três linhas de
cada tabela foi o melhor formato, com 67,0\% de acurácia de execução no
conjunto de desenvolvimento do Spider, à frente de listas de colunas e de
documentação em estilo de API \cite{rajkumar2022evaluating}. O DAIL-SQL
comparou cinco representações de pergunta e esquema em vários modelos e
mediu o efeito de detalhes do prompt: incluir as chaves estrangeiras
melhorou a acurácia de execução em 0,6 a 2,9 pontos percentuais na
maioria dos casos, e retirar a instrução de não explicar a resposta a
reduziu em 1,3 a 2,4 pontos \cite{gao2024dail}. Pequenas regras sobre o
formato da resposta, portanto, já têm efeito mensurável.

Outros trabalhos mudam o processo, e não só o prompt. O DIN-SQL decompõe
a geração em etapas, com seleção das tabelas, classificação da pergunta
por complexidade, geração e autocorreção, e chegou a 85,3\% de acurácia
de execução no teste do Spider e 55,9\% no do BIRD
\cite{pourreza2023din}. O CHESS usa agentes para recuperar do catálogo as
descrições e os valores relevantes e para podar esquemas grandes; a poda
deu cerca de 2 pontos percentuais a mais com cinco vezes menos
\textit{tokens}, e o sistema chegou a 71,10\% no teste do BIRD
\cite{talaei2024chess}.

Uma terceira linha acrescenta conhecimento ao prompt e mede o ganho. No
BIRD, fornecer a evidência de cada pergunta elevou a acurácia de execução
do GPT-4 no teste de 34,88\% para 54,89\% \cite{li2023bird}. Em uma
ablação incremental com um modelo de contexto longo, a mesma evidência
levou a acurácia de 43,87\% para 60,23\%, e os acréscimos seguintes,
valores de exemplo das colunas e autocorreção, levaram a 61,99\% e
64,80\%; a documentação de SQL quase não alterou o resultado
\cite{chung2025longcontext}. Variando apenas as descrições de colunas no
BIRD, descrições geradas pelo Qwen2 elevaram a acurácia do GPT-4o de
0,301 para 0,368 e a do Mixtral 8x22B de 0,175 para 0,296; já as
descrições geradas pelo GPT-4o reduziram a do Qwen2-72B de 0,289 para
0,260 \cite{wretblad2024synthetic}. Descrições ajudam, mas uma descrição
pior pode atrapalhar.

O estudo mais próximo deste trabalho é uma ablação sobre um banco de
telemetria da Intel, com 25 tabelas, 30 perguntas e o modelo proprietário
GPT-4.1 \cite{bouattour2026prompting}. Partindo de 43,3\% de acurácia, as
descrições de tabelas e colunas tiradas do dicionário de dados levaram a
56,7\%, os exemplos resolvidos a 53,3\%, e a combinação dos dois a
83,3\%; com validação das consultas, a configuração completa chegou a
90,0\%. A configuração completa também foi a mais cara, com custo por
consulta cerca de 3,7 vezes o da configuração de partida. Os autores
associam o peso das descrições aos nomes abreviados e à terminologia
específica do esquema de telemetria e indicam como trabalho futuro
testar modelos abertos menores.

Essas ablações separam formatos, exemplos e descrições, mas não separam
as descrições de colunas das regras de consulta válidas para a base
inteira: grão, junções com integridade imperfeita e formato da resposta.
A evidência do BIRD mistura os dois tipos e é escrita para cada pergunta;
o dicionário de dados da Intel corresponde ao nível de catálogo. Este
trabalho mantém fixos o modelo, as instruções e o formato da resposta, e
varia apenas o conteúdo do contexto, em três níveis aditivos
(\autoref{sec:condicoes}).

% -----------------------------------------------------------------------------
\section{Camadas semânticas e modelos de linguagem}
\label{sec:industria}

Ferramentas de engenharia e de análise de dados mantêm uma camada
semântica: um conjunto de definições de métricas, grãos, dimensões
temporais e relacionamentos, mantido fora das tabelas e usado para gerar
as consultas \cite{jones2024semanticlayer}. São exemplos o dbt Semantic
Layer \cite{dbt_semanticlayer}, as \textit{metric views} do Databricks
Unity Catalog \cite{databricks_metricviews}, o LookML do Looker
\cite{looker_lookml} e o Cube \cite{cube_datamodel}, descritos na
\autoref{sec:camada-semantica}. A correspondência entre essas ferramentas
e o contexto semântico deste trabalho está no \autoref{qua:industriais}.

Os fornecedores e a comunidade passaram a medir o efeito dessas
definições sobre agentes de linguagem. Um relato publicado na comunidade
Databricks acrescentou camadas cumulativas de contexto ao agente da
plataforma, com 30 perguntas e duas execuções por camada, e registrou 40\%
de acurácia com apenas as tabelas, 47\% com as descrições do catálogo e
97\% com as \textit{metric views}; exemplos de SQL e instruções em texto,
acrescentados depois, levaram a 100\% e 97\% \cite{databricks2026genie}.
O relato não informa o modelo usado pelo agente nem traz teste
estatístico.

A dbt Labs comparou o modelo escrevendo SQL diretamente com o modelo
consultando a camada semântica, que gera o SQL
\cite{ganz2023semanticlayer,ganz2026benchmark}. Na versão de 2026, sobre
uma base de seguros de 15 tabelas, com 11 perguntas executadas 20 vezes
cada, a camada semântica acertou 98,2\% e 100\% das respostas com dois
modelos proprietários, contra 90,0\% e 84,1\% do SQL escrito pelo modelo
\cite{ganz2026benchmark}. Fora do escopo das métricas definidas, a camada
semântica não responde, em vez de errar. O próprio relato registra que a
empresa vende o produto avaliado e que o estudo usa uma única base. Sobre
a mesma base de seguros, um estudo anterior mostrou que o GPT-4 acertava
16\% das perguntas escrevendo SQL sobre o esquema e 54\% consultando uma
representação do banco em grafo de conhecimento, com ontologia e
mapeamentos \cite{sequeda2023benchmark}.

Trabalhos acadêmicos recentes também anexam definições semânticas ao
prompt, sem compilador. Um estudo pareado com 100 perguntas sobre uma
base de varejo comparou o esquema sozinho com o esquema acrescido de um
documento semântico de 4~KB, escrito à mão, com medidas, convenções e
regras de desambiguação \cite{rumiantsau2026semantic}. Com três modelos
proprietários, a acurácia passou de 45,5--50,5\% para 67,7--68,7\%, um
ganho de 17 a 23 pontos percentuais, significativo em todos os modelos; as
diferenças entre os modelos, dentro de cada condição, não foram
significativas. Outro trabalho propõe definições governadas de métricas,
junções e grão, com validação das consultas geradas, para reduzir
respostas inventadas em análises corporativas \cite{pillai2026ground}.

Esses resultados apontam na mesma direção, mas cada um deixa uma
pergunta. O relato da Databricks e os da dbt Labs medem produtos
comerciais, em que o conteúdo do contexto e a plataforma não se separam.
Os estudos de grafo de conhecimento e de camada semântica compilada
mudam a forma de consultar, e não só o contexto. E o documento semântico
único não separa as descrições das regras. Nenhum deles usa modelos de
pesos abertos para isolar, em níveis aditivos, a contribuição do catálogo
e a das regras.

% -----------------------------------------------------------------------------
\section{Governança de dados para agentes}
\label{sec:governanca}

Um agente que gera e executa SQL sobre dados de uma organização levanta
duas questões além da acurácia: o que é enviado ao provedor do modelo e o
que a consulta gerada pode fazer no banco. Ataques de injeção por meio do
prompt conseguem levar aplicações integradas a modelos de linguagem a
executar consultas SQL não pretendidas \cite{pedro2025prompt}. O estudo
que caracterizou esses ataques os reproduziu em sete modelos de linguagem
e propôs quatro defesas: restringir as permissões do usuário de banco
usado pela aplicação, de modo que só leituras sejam possíveis; reescrever
as consultas geradas para limitar os registros visíveis a cada usuário;
usar um segundo modelo, sem acesso ao banco, para inspecionar os
resultados em busca de injeções; e carregar previamente no prompt os
dados do usuário, para que a aplicação não precise consultar dados de
outros \cite{pedro2025prompt}. A lista de riscos da OWASP para aplicações
com modelos de linguagem inclui, entre outros, a injeção de prompt, a
divulgação de informação sensível, o tratamento inadequado da saída do
modelo e o excesso de autonomia concedida ao agente
\cite{owasp2025llm}.

Este trabalho adota a primeira dessas defesas e uma restrição de
privacidade: a consulta gerada passa por um validador que só aceita
leitura e é executada por um usuário sem permissão de escrita
(\autoref{sec:agente}), e nenhum valor de linha é enviado ao provedor
(\autoref{sec:privacidade}). A camada semântica também tem papel de
governança: ela fixa o que a organização entende por cada métrica e por
cada relacionamento, em vez de deixar essa definição ao modelo.

% -----------------------------------------------------------------------------
\section{Síntese e lacuna}
\label{sec:lacuna}

O \autoref{qua:relacionados} resume os trabalhos mais próximos segundo
quatro critérios do desenho deste trabalho: a base, a variação de
contexto, os modelos e a dependência de plataforma proprietária.

\begin{quadro}[htb]
  \centering
  \caption{Comparação com trabalhos relacionados}
  \label{qua:relacionados}
  \footnotesize
  \begin{tabularx}{\textwidth}{|>{\raggedright\arraybackslash}p{2.9cm}|*{4}{>{\raggedright\arraybackslash}X|}}
    \hline
    \textbf{Trabalho} & \textbf{Base de dados} & \textbf{Variação de
      contexto} & \textbf{Modelos} & \textbf{Plataforma} \\ \hline
    \citeonline{li2023bird} & 95 bancos públicos, dados reais, em inglês
      & evidência por pergunta, com e sem & abertos e proprietários &
      nenhuma \\ \hline
    \citeonline{wretblad2024synthetic} & BIRD & só as descrições de
      colunas, de várias origens & abertos e proprietários & nenhuma \\
      \hline
    \citeonline{bouattour2026prompting} & telemetria da Intel, 25 tabelas
      & descrições, exemplos e validação, em ablação & GPT-4.1
      (proprietário) & nenhuma \\ \hline
    \citeonline{rumiantsau2026semantic} & base de exemplo de varejo &
      esquema com e sem documento semântico único & três modelos
      proprietários & nenhuma \\ \hline
    \citeonline{sequeda2023benchmark} & esquema de seguros & SQL direto
      ou consulta a grafo de conhecimento & GPT-4 (proprietário) & grafo
      de conhecimento \\ \hline
    \citeonline{ganz2026benchmark} & esquema de seguros, 15 tabelas & SQL
      direto ou consulta à camada semântica & modelos proprietários & dbt
      Semantic Layer \\ \hline
    \citeonline{databricks2026genie} & base de exemplo da plataforma &
      camadas cumulativas: tabelas, descrições, \textit{metric views},
      exemplos, instruções & não informado & proprietária \\ \hline
    Este trabalho & base comercial real descaracterizada, 11 tabelas, em
      português, com referências órfãs & três níveis aditivos: técnico,
      catálogo, semântico & dois modelos de pesos abertos de famílias
      distintas & nenhuma: arquivos de texto, PostgreSQL e API gratuita
      \\ \hline
  \end{tabularx}
  \fonte{Elaborado pelo autor a partir dos trabalhos citados. Linha deste
    trabalho: \arq{docs/briefing.md} e \arq{docs/decisions.md} (D-001,
    D-004, D-008).}
\end{quadro}

Nenhum dos trabalhos responde, ao mesmo tempo, se o efeito da camada
semântica aparece quando ela é apenas texto enviado ao modelo, sobre uma
base comercial real, com modelos de pesos abertos, e qual parte do
contexto o produz. Os \textit{benchmarks} medem modelos sobre bases
públicas; as ablações de contexto separam descrições, exemplos e formatos,
mas não descrições e regras; os estudos de camada semântica usam modelos
proprietários, um documento único ou uma forma de consulta diferente; e o
relato com níveis cumulativos vem de uma plataforma proprietária. Este
trabalho ocupa essa interseção (\autoref{sec:objetivos}).
